\documentclass[12pt]{article}

\usepackage[margin=1in]{geometry}
\usepackage{amsmath}
\usepackage{array}
\usepackage{tikz}

\setlength{\parindent}{0pt}
\setlength{\parskip}{0.7em}

\newcommand{\CorrelationScatter}[2]{%
	\begin{tikzpicture}[x=0.32cm,y=0.32cm]
		\draw[->] (0,0) -- (10.5,0);
		\draw[->] (0,0) -- (0,10.5);
		\node[font=\scriptsize, below] at (10.5,0) {$x$};
		\node[font=\scriptsize, left] at (0,10.5) {$y$};
		\foreach \x/\y in {#2}{
			\fill (\x,\y) circle (1.2pt);
		}
		\node[font=\small] at (5.25,11.1) {#1};
	\end{tikzpicture}%
}

\begin{document}
	
	\section{Explanatory material}
	
	\subsection{Calculating the linear correlation coefficient $r$}
	\textbf{Cycle 2 -- April 7th}
	
	\subsubsection*{What is the linear correlation coefficient?}
	
	The linear correlation coefficient ($r$) measures the strength and direction of the relationship between two variables.
	
	\[
	\begin{aligned}
		r=+1 &\rightarrow \text{perfect positive relationship} \\
		r=0  &\rightarrow \text{no linear relationship} \\
		r=-1 &\rightarrow \text{perfect negative relationship}
	\end{aligned}
	\]
	
	Examples:
	\begin{itemize}
		\item Positive correlation
		\item Negative correlation
		\item No correlation
	\end{itemize}
	
	\subsubsection*{Representative scatterplots}
	
	\begin{center}
		\begin{tabular}{ccc}
			\CorrelationScatter{$r=-1$}{
				0.250/8.958, 0.750/8.542, 1.250/8.125, 1.750/7.708, 2.250/7.292, 2.750/6.875, 3.250/6.458, 3.750/6.042, 4.250/5.625, 4.750/5.208, 5.250/4.792, 5.750/4.375, 6.250/3.958, 6.750/3.542, 7.250/3.125, 7.750/2.708, 8.250/2.292, 8.750/1.875, 9.250/1.458, 9.750/1.042
			}
			&
			\CorrelationScatter{$r=-0.8$}{
				0.250/9.042, 0.750/6.458, 1.250/5.625, 1.750/7.042, 2.250/6.958, 2.750/8.375, 3.250/3.792, 3.750/7.458, 4.250/7.875, 4.750/4.292, 5.250/6.208, 5.750/5.625, 6.250/4.542, 6.750/2.708, 7.250/1.375, 7.750/2.542, 8.250/2.458, 8.750/4.625, 9.250/2.792, 9.750/0.208
			}
			&
			\CorrelationScatter{$r=-0.5$}{
				0.250/8.100, 0.750/4.795, 1.250/3.906, 1.750/6.123, 2.250/6.268, 2.750/8.484, 3.250/2.420, 3.750/7.740, 4.250/8.576, 4.750/3.892, 5.250/6.798, 5.750/6.253, 6.250/5.020, 6.750/2.750, 7.250/1.172, 7.750/3.042, 8.250/3.188, 8.750/6.438, 9.250/4.169, 9.750/0.865
			}
			\\[1em]
			\CorrelationScatter{$r=-0.3$}{
				0.250/7.579, 0.750/3.876, 1.250/2.956, 1.750/5.614, 2.250/5.886, 2.750/8.544, 3.250/1.661, 3.750/7.896, 4.250/8.964, 4.750/3.671, 5.250/7.124, 5.750/6.601, 6.250/5.284, 6.750/2.774, 7.250/1.059, 7.750/3.319, 8.250/3.591, 8.750/7.441, 9.250/4.931, 9.750/1.229
			}
			&
			\CorrelationScatter{$r=0$}{
				0.250/6.458, 0.750/2.708, 1.250/1.875, 1.750/4.792, 2.250/5.208, 2.750/8.125, 3.250/1.042, 3.750/7.708, 4.250/8.958, 4.750/3.542, 5.250/7.292, 5.750/6.875, 6.250/5.625, 6.750/3.125, 7.250/1.458, 7.750/3.958, 8.250/4.375, 8.750/8.542, 9.250/6.042, 9.750/2.292
			}
			&
			\CorrelationScatter{$r=0.3$}{
				0.250/5.204, 0.750/1.751, 1.250/1.081, 1.750/3.989, 2.250/4.511, 2.750/7.419, 3.250/0.786, 3.750/7.271, 4.250/8.589, 4.750/3.546, 5.250/7.249, 5.750/6.976, 6.250/5.909, 6.750/3.649, 7.250/2.184, 7.750/4.694, 8.250/5.216, 8.750/9.316, 9.250/7.056, 9.750/3.604
			}
			\\[1em]
			\CorrelationScatter{$r=0.5$}{
				0.250/4.315, 0.750/1.410, 1.250/0.919, 1.750/3.533, 2.250/4.077, 2.750/6.691, 3.250/1.025, 3.750/6.744, 4.250/7.979, 4.750/3.693, 5.250/6.997, 5.750/6.851, 6.250/6.015, 6.750/4.145, 7.250/2.964, 7.750/5.233, 8.250/5.777, 8.750/9.426, 9.250/7.555, 9.750/4.650
			}
			&
			\CorrelationScatter{$r=0.8$}{
				0.250/2.708, 0.750/0.792, 1.250/0.625, 1.750/2.708, 2.250/3.292, 2.750/5.375, 3.250/1.458, 3.750/5.792, 4.250/6.875, 4.750/3.958, 5.250/6.542, 5.750/6.625, 6.250/6.208, 6.750/5.042, 7.250/4.375, 7.750/6.208, 8.250/6.792, 8.750/9.625, 9.250/8.458, 9.750/6.542
			}
			&
			\CorrelationScatter{$r=1$}{
				0.250/1.042, 0.750/1.458, 1.250/1.875, 1.750/2.292, 2.250/2.708, 2.750/3.125, 3.250/3.542, 3.750/3.958, 4.250/4.375, 4.750/4.792, 5.250/5.208, 5.750/5.625, 6.250/6.042, 6.750/6.458, 7.250/6.875, 7.750/7.292, 8.250/7.708, 8.750/8.125, 9.250/8.542, 9.750/8.958
			}
		\end{tabular}
	\end{center}
	
	\subsubsection*{Linear Correlation Coefficient ($r$)}
	
	\[
	r=
	\frac{n\sum xy-(\sum x)(\sum y)}
	{\sqrt{\left[n\sum x^2-(\sum x)^2\right]\left[n\sum y^2-(\sum y)^2\right]}}
	\]
	
	\textbf{Meaning of each part:}
	\begin{itemize}
		\item $r$: The correlation coefficient. It ranges from $-1$ to $+1$.
		\item $n$: Number of data points (observations).
		\item $\sum x$: Sum of all $x$-values.
		\item $\sum y$: Sum of all $y$-values.
		\item $\sum xy$: Sum of the products of each $x$ and $y$ pair.
		\item $\sum x^2$: Sum of the squares of all $x$-values.
		\item $\sum y^2$: Sum of the squares of all $y$-values.
	\end{itemize}
	
	\textbf{Interpretation of $r$:}
	\begin{itemize}
		\item $-1$: Perfect negative correlation
		\item $0$: No linear correlation
		\item $+1$: Perfect positive correlation
	\end{itemize}
	
	The closer $r$ is to $-1$ or $+1$, the stronger the linear relationship.
	
	\subsubsection*{Properties of the Linear Correlation Coefficient $r$}
	
	\begin{enumerate}
		\item The value of $r$ is always between $-1$ and $1$ inclusive. That is,
		\[
		-1 \leq r \leq 1
		\]
		
		\item If all values of either variable are converted to a different scale, the value of $r$ does not change.
		
		\item The value of $r$ is not affected by the choice of $x$ or $y$. Interchange all $x$- and $y$-values and the value of $r$ will not change.
		
		\item $r$ measures the strength of a linear relationship. It is not designed to measure the strength of a relationship that is not linear.
		
		\item $r$ is very sensitive to outliers in the sense that a single outlier can dramatically affect its value.
	\end{enumerate}
	
	\subsubsection*{Why is it useful in real life?}
	
	Understanding correlation helps you:
	\begin{itemize}
		\item Make better decisions using data
		\item Understand how variables are connected in real situations
		\item Predict possible trends in economics and daily life
		\item Analyse information in studies, sport, and finance
	\end{itemize}
	
	\subsection{Calculating the linear regression equation}
	\textbf{Cycle 3 -- April 9th}
	
	\subsubsection*{The linear regression equation using the least squares method}
	
	\[
	\hat{y}=a+bx
	\]
	
	Where:
	\begin{itemize}
		\item $\hat{y}$ is the predicted value,
		\item $a$ is the intercept,
		\item $b$ is the slope,
		\item $x$ is the independent variable.
	\end{itemize}
	
	\subsubsection*{Linear regression}
	
	To calculate the slope $b$:
	\[
	b=\frac{n\sum(xy)-\sum x \sum y}{n\sum x^2-(\sum x)^2}
	\]
	
	To calculate the intercept $a$:
	\[
	a=\frac{\sum y-b\sum x}{n}
	\]
	
	\subsubsection*{Relationship between the linear correlation coefficient and the regression slope}
	
	The linear correlation coefficient $r$ and the regression slope $b$ are related, but they are not the same quantity.
	
	The linear correlation coefficient $r$ measures the strength and direction of a linear relationship between two variables. The regression slope $b$ in
	\[
	\hat{y}=a+bx
	\]
	measures how much the predicted value of $y$ changes when $x$ increases by $1$ unit.
	
	So, $r$ is a measure of association, while $b$ is a measure of change.
	
	The formula for the linear correlation coefficient is
	\[
	r=
	\frac{n\sum xy-(\sum x)(\sum y)}
	{\sqrt{\left[n\sum x^2-(\sum x)^2\right]\left[n\sum y^2-(\sum y)^2\right]}}.
	\]
	
	The formula for the regression slope is
	\[
	b=
	\frac{n\sum xy-(\sum x)(\sum y)}
	{n\sum x^2-(\sum x)^2}.
	\]
	
	Notice that the numerator is the same in both formulas:
	\[
	n\sum xy-(\sum x)(\sum y).
	\]
	However, the denominator is different. This is one reason why $r$ and $b$ are generally not equal.
	
	The exact relationship between them is
	\[
	b=r\frac{s_y}{s_x},
	\]
	where $s_y$ is the standard deviation of $y$ and $s_x$ is the standard deviation of $x$.
	
	This formula shows that $b$ depends not only on the direction and strength of the linear relationship, but also on the relative spread of the two variables. For this reason, $r$ and $b$ may sometimes look similar in value, but they are usually different.
	
	It is also important to interpret them correctly:
	\begin{itemize}
		\item $r$ is always between $-1$ and $1$.
		\item $r$ has no units.
		\item $b$ is not restricted to the interval $[-1,1]$.
		\item $b$ has units, specifically units of $y$ per unit of $x$.
	\end{itemize}
	
	In a specific example with nonzero correlation, if $r=b$, then
	\[
	\frac{s_y}{s_x}=1,
	\]
	so the spread of $x$ and $y$ is the same. This is a special case, not a general rule.
	
	A useful memory aid is:
	\begin{itemize}
		\item $r$ answers: ``How strong is the linear relationship?''
		\item $b$ answers: ``How much does the predicted $y$ change when $x$ increases by $1$?''
	\end{itemize}
	
	For terminology, $r$ should be called the linear correlation coefficient, and $b$ should be called the regression slope or slope coefficient.
	
	\section{Activities and solutions}
	
	\subsection{Task one}
	
	Complete the following table and calculate the Pearson product moment correlation coefficient. Draw the corresponding scatter plot.
	
	\begin{center}
		\renewcommand{\arraystretch}{2}
		\begin{tabular}{|c|c|c|c|c|}
			\hline
			$(x)$ & $(y)$ & $x^2$ & $y^2$ & $xy$ \\ \hline
			1 & 2 & 1 & 4 & 2 \\ \hline
			2 & 3 & 4 & 9 & 6 \\ \hline
			3 & 5 & 9 & 25 & 15 \\ \hline
			4 & 4 & 16 & 16 & 16 \\ \hline
			5 & 6 & 25 & 36 & 30 \\ \hline
			$\sum 15$ & $\sum 20$ & $\sum 55$ & $\sum 90$ & $\sum 69$ \\ \hline
		\end{tabular}
	\end{center}
	
	\subsubsection*{Solution}
	
	Using the values in the table, we have
	\[
	n=5,\qquad \sum x=15,\qquad \sum y=20,\qquad \sum x^2=55,\qquad \sum y^2=90,\qquad \sum xy=69.
	\]
	
	Substitute into the Pearson product moment correlation coefficient formula:
	\[
	r=
	\frac{n\sum xy-(\sum x)(\sum y)}
	{\sqrt{\left[n\sum x^2-(\sum x)^2\right]\left[n\sum y^2-(\sum y)^2\right]}}.
	\]
	
	So,
	\[
	r=
	\frac{5(69)-(15)(20)}
	{\sqrt{\left[5(55)-15^2\right]\left[5(90)-20^2\right]}}.
	\]
	
	Now simplify:
	\[
	r=
	\frac{345-300}
	{\sqrt{(275-225)(450-400)}}
	=
	\frac{45}{\sqrt{50\cdot 50}}
	=
	\frac{45}{50}
	=0.9.
	\]
	
	Therefore, the Pearson product moment correlation coefficient is
	\[
	r=0.9.
	\]
	
	This indicates a strong positive linear correlation between \(x\) and \(y\).
	
	The corresponding scatter plot is shown below.
	
	\begin{center}
		\begin{tikzpicture}[x=1.2cm,y=1.0cm]
			\draw[->] (0,0) -- (6.3,0) node[right] {$x$};
			\draw[->] (0,0) -- (0,6.8) node[above] {$y$};
			
			\foreach \x in {1,2,3,4,5}
			\draw (\x,0.08) -- (\x,-0.08) node[below] {\x};
			
			\foreach \y in {1,2,3,4,5,6}
			\draw (0.08,\y) -- (-0.08,\y) node[left] {\y};
			
			\foreach \x/\y in {1/2,2/3,3/5,4/4,5/6}
			\fill (\x,\y) circle (2pt);
		\end{tikzpicture}
	\end{center}
	
	Now calculate the linear regression equation \(\hat{y}=a+bx\).
	
	First, compute the slope:
	\[
	b=\frac{n\sum xy-(\sum x)(\sum y)}{n\sum x^2-(\sum x)^2}.
	\]
	
	Substitute the values:
	\[
	b=\frac{5(69)-(15)(20)}{5(55)-15^2}
	=\frac{345-300}{275-225}
	=\frac{45}{50}
	=0.9.
	\]
	
	Now compute the intercept:
	\[
	a=\frac{\sum y-b\sum x}{n}.
	\]
	
	Substitute the values:
	\[
	a=\frac{20-(0.9)(15)}{5}
	=\frac{20-13.5}{5}
	=\frac{6.5}{5}
	=1.3.
	\]
	
	Therefore, the linear regression equation is
	\[
	\hat{y}=1.3+0.9x.
	\]
	
	The regression line is shown below.
	
	\begin{center}
		\begin{tikzpicture}[x=1.2cm,y=1.0cm]
			\draw[->] (0,0) -- (6.3,0) node[right] {$x$};
			\draw[->] (0,0) -- (0,6.8) node[above] {$y$};
			
			\foreach \x in {1,2,3,4,5}
			\draw (\x,0.08) -- (\x,-0.08) node[below] {\x};
			
			\foreach \y in {1,2,3,4,5,6}
			\draw (0.08,\y) -- (-0.08,\y) node[left] {\y};
			
			\draw[thick] (0,1.3) -- (5.5,6.25);
			
			\draw[->] (0,0) -- (6.3,0) node[right] {$x$};
			\draw[->] (0,0) -- (0,6.8) node[above] {$y$};
			
			\foreach \x in {1,2,3,4,5}
			\draw (\x,0.08) -- (\x,-0.08) node[below] {\x};
			
			\foreach \y in {1,2,3,4,5,6}
			\draw (0.08,\y) -- (-0.08,\y) node[left] {\y};
			
			\foreach \x/\y in {1/2,2/3,3/5,4/4,5/6}
			\fill (\x,\y) circle (2pt);
		\end{tikzpicture}
	\end{center}
	
	
	\subsection{Warm up}
	
	Calculate the question mark value.
	
	\[
	\begin{aligned}
		\text{accordion} + \text{accordion} + \text{accordion} &= 36 \\
		\text{guitar} \times \text{guitar} + \text{accordion} &= 31 \\
		\text{pair of guitars} + \text{pair of guitars} \times \text{pair of banjos} &= 190 \\
		\text{accordion} + \text{banjo} \times \text{guitar} &= ?
	\end{aligned}
	\]
	
	\subsubsection*{Solution}
	
	Interpreting each double-instrument icon as twice the corresponding single-instrument value:
	
	\[
	3(\text{accordion})=36
	\]
	so
	\[
	\text{accordion}=12.
	\]
	
	From the second line,
	\[
	(\text{guitar})^2+12=31
	\]
	so
	\[
	(\text{guitar})^2=19
	\]
	and therefore
	\[
	\text{guitar}=\sqrt{19}.
	\]
	
	From the third line,
	\[
	2(\text{guitar})+2(\text{guitar})\cdot 2(\text{banjo})=190.
	\]
	Substituting \(\text{guitar}=\sqrt{19}\),
	\[
	2\sqrt{19}+4\sqrt{19}\,(\text{banjo})=190,
	\]
	hence
	\[
	4\sqrt{19}\,(\text{banjo})=190-2\sqrt{19},
	\]
	\[
	\text{banjo}=\frac{190-2\sqrt{19}}{4\sqrt{19}}
	=\frac{95-\sqrt{19}}{2\sqrt{19}}.
	\]
	
	Now evaluate the final expression:
	\[
	\text{accordion}+\text{banjo}\cdot\text{guitar}
	=12+\left(\frac{95-\sqrt{19}}{2\sqrt{19}}\right)\sqrt{19}.
	\]
	So
	\[
	?=12+\frac{95-\sqrt{19}}{2}
	=\frac{119-\sqrt{19}}{2}.
	\]
	
	Therefore,
	\[
	?=\frac{119-\sqrt{19}}{2}\approx 57.32.
	\]
	
	\subsection{Activity 1}
	
	\begin{enumerate}
		\item Calculate the linear regression equation.
		\item Create the scatter plot with the linear regression equation.
	\end{enumerate}
	
	\begin{center}
		\small
		\setlength{\tabcolsep}{6pt}
		\begin{tabular}{|c|c|c|c|c|c|c|c|c|c|}
			\hline
			Test 1 & 10 & 18 & 13 & 6 & 8 & 5 & 12 & 15 & 15 \\ \hline
			Test 2 & 12 & 20 & 11 & 9 & 6 & 6 & 12 & 13 & 17 \\ \hline
		\end{tabular}
	\end{center}
	
	\subsubsection*{Solution}
	
	Take Test 1 as \(x\) and Test 2 as \(y\). The data pairs are
	\[
	(10,12),\ (18,20),\ (13,11),\ (6,9),\ (8,6),\ (5,6),\ (12,12),\ (15,13),\ (15,17).
	\]
	
	Since there are \(9\) pairs,
	\[
	n=9.
	\]
	
	Now calculate the required sums directly from the table.
	
	\[
	\sum x = 10+18+13+6+8+5+12+15+15 = 102
	\]
	
	\[
	\sum y = 12+20+11+9+6+6+12+13+17 = 106
	\]
	
	\[
	\sum x^2 = 10^2+18^2+13^2+6^2+8^2+5^2+12^2+15^2+15^2
	\]
	\[
	=100+324+169+36+64+25+144+225+225 = 1312
	\]
	
	\[
	\sum xy = (10)(12)+(18)(20)+(13)(11)+(6)(9)+(8)(6)+(5)(6)+(12)(12)+(15)(13)+(15)(17)
	\]
	\[
	=120+360+143+54+48+30+144+195+255 = 1349
	\]
	
	Substitute these values into the slope formula:
	\[
	b=\frac{n\sum(xy)-\sum x\sum y}{n\sum x^2-(\sum x)^2}
	\]
	
	\[
	b=\frac{9(1349)-(102)(106)}{9(1312)-102^2}
	\]
	
	\[
	b=\frac{12141-10812}{11808-10404}
	\]
	
	\[
	b=\frac{1329}{1404}
	\]
	
	\[
	b=\frac{443}{468}\approx 0.947
	\]
	
	Now calculate the intercept:
	\[
	a=\frac{\sum y-b\sum x}{n}
	\]
	
	\[
	a=\frac{106-\left(\frac{443}{468}\right)(102)}{9}
	\]
	
	\[
	a=\frac{106-\frac{15062}{468}}{9}
	\]
	
	\[
	a=\frac{\frac{49608}{468}-\frac{15062}{468}}{9}
	=\frac{\frac{34546}{468}}{9}
	\]
	
	\[
	a=\frac{34546}{4212}
	=\frac{17273}{2106}
	=\frac{737}{702}
	\approx 1.050
	\]
	
	Therefore, the linear regression equation is
	\[
	\hat{y}=a+bx
	\]
	
	\[
	\hat{y}=1.050+0.947x
	\]
	
	For the scatter plot, plot the nine points
	\[
	(10,12),\ (18,20),\ (13,11),\ (6,9),\ (8,6),\ (5,6),\ (12,12),\ (15,13),\ (15,17)
	\]
	and then draw the regression line
	\[
	\hat{y}=1.050+0.947x.
	\]
	
	The scatter plot with the regression line is shown below.
	
	\begin{center}
		\begin{tikzpicture}[x=0.42cm,y=0.35cm]
			\draw[->] (0,0) -- (15.5,0) node[right] {$x$};
			\draw[->] (0,0) -- (0,16.5) node[above] {$y$};
			
			\foreach \x/\y in {6/7,14/15,9/6,2/4,4/1,1/1,8/7,11/8,11/12}{
				\fill (\x,\y) circle (1.4pt);
			}
			
			\draw[thick] (1,0.783) -- (14,13.089);
			
			\node[font=\small] at (7.75,17.2) {Scatter plot and regression line};
		\end{tikzpicture}
	\end{center}
	
	Here the plotted coordinates are the translated data points \((x-4,y-5)\), so the relative positions are preserved and the line shown corresponds to the same regression relationship.
	
	\subsection{Activity 2 -- Data set 2}
	
	\begin{enumerate}
		\item Calculate the linear regression equation.
		\item Create the scatter plot with the linear regression equation.
	\end{enumerate}
	
	\begin{center}
		\small
		\setlength{\tabcolsep}{6pt}
		\begin{tabular}{|>{\raggedright\arraybackslash}m{4.5cm}|c|c|c|c|c|c|}
			\hline
			\shortstack[l]{Minimum temperature $(x)$} & 17.7 & 19.8 & 23.3 & 22.4 & 22.0 & 22.0 \\ \hline
			\shortstack[l]{Maximum temperature $(y)$} & 29.4 & 34.0 & 34.5 & 35.0 & 36.9 & 36.4 \\ \hline
		\end{tabular}
	\end{center}
	
	The table above shows the maximum and minimum temperatures (in \(^{\circ}\mathrm{C}\)) during a hot week in Melbourne.
	
	\subsubsection*{Solution}
	
	For this data set,
	\[
	n=6,\qquad \sum x=127.2,\qquad \sum y=206.2,
	\]
	\[
	\sum x^2=2717.98,\qquad \sum xy=4394.03.
	\]
	
	The slope is
	\[
	b=\frac{6(4394.03)-127.2(206.2)}{6(2717.98)-(127.2)^2}
	\approx 1.059.
	\]
	
	The intercept is
	\[
	a=\frac{206.2-1.058575445(127.2)}{6}\approx 11.925.
	\]
	
	Therefore, the linear regression equation is
	\[
	\hat{y}=11.925+1.059x.
	\]
	
	The scatter plot with the regression line is shown below.
	
	\begin{center}
		\begin{tikzpicture}[x=0.8cm,y=0.35cm]
			\draw[->] (0,0) -- (8.5,0) node[right] {$x$};
			\draw[->] (0,0) -- (0,10.5) node[above] {$y$};
			
			\foreach \x/\y in {1.7/1.4,3.8/6.0,7.3/6.5,6.4/7.0,6.0/8.9,6.0/8.4}{
				\fill (\x,\y) circle (1.4pt);
			}
			
			\draw[thick] (1.7,2.662) -- (7.3,8.590);
			
			\node[font=\small] at (4.25,11.2) {Scatter plot and regression line};
		\end{tikzpicture}
	\end{center}
	
	Here the plotted coordinates are the translated data points \((x-16,y-28)\), so the relative positions are preserved and the line shown corresponds to the same regression relationship.
	
	\subsection{Activity 2 -- Data set 3}
	
	\begin{enumerate}
		\item Calculate the linear regression equation.
		\item Create the scatter plot with the linear regression equation.
	\end{enumerate}
	
	\begin{center}
		\small
		\setlength{\tabcolsep}{5pt}
		\begin{tabular}{|>{\raggedright\arraybackslash}m{3cm}|c|c|c|c|c|c|c|c|c|c|c|}
			\hline
			\shortstack[l]{Balls faced} & 29 & 16 & 19 & 62 & 13 & 40 & 16 & 9 & 28 & 26 & 6 \\ \hline
			\shortstack[l]{Runs scored} & 27 & 8 & 21 & 47 & 3 & 15 & 13 & 2 & 15 & 10 & 2 \\ \hline
		\end{tabular}
	\end{center}
	
	The table above shows the number of runs scored and the number of balls faced by batsmen in a 1-day international cricket match.
	
	\subsubsection*{Solution}
	
	For this data set,
	\[
	n=11,\qquad \sum x=264,\qquad \sum y=163,
	\]
	\[
	\sum x^2=8904,\qquad \sum xy=5781.
	\]
	
	The slope is
	\[
	b=\frac{11(5781)-264(163)}{11(8904)-264^2}
	=\frac{1946}{2674}\approx 0.728.
	\]
	
	The intercept is
	\[
	a=\frac{163-0.727803738(264)}{11}\approx -2.649.
	\]
	
	Therefore, the linear regression equation is
	\[
	\hat{y}=-2.649+0.728x.
	\]
	
	The scatter plot with the regression line is shown below.
	
	\begin{center}
		\begin{tikzpicture}[x=0.55cm,y=0.45cm]
			\draw[->] (0,0) -- (13.2,0) node[right] {$x$};
			\draw[->] (0,0) -- (0,10.5) node[above] {$y$};
			
			\foreach \x/\y in {5.8/5.4,3.2/1.6,3.8/4.2,12.4/9.4,2.6/0.6,8.0/3.0,3.2/2.6,1.8/0.4,5.6/3.0,5.2/2.0,1.2/0.4}{
				\fill (\x,\y) circle (1.4pt);
			}
			
			\draw[thick] (1.2,0.344) -- (12.4,8.495);
			
			\node[font=\small] at (6.6,11.2) {Scatter plot and regression line};
		\end{tikzpicture}
	\end{center}
	
	Here the plotted coordinates are the scaled data points \((x/5,y/5)\), so the relative positions are preserved and the line shown corresponds to the same regression relationship.
	
\end{document}	